\scanpage{89}
\mspage{82}
% archive sheet number: 88
\[
\operatorname{card}I'_\alpha\leq
\underbrace{(\operatorname{card}\alpha)\pi}_{\leq\pi.\pi=\pi^2=\pi}
\qquad\text{puisque les }\operatorname{card}I'_\beta\leq\pi,
\]
on gagne encore.

On trouve enfin, par les $=$ arguments
\[
\operatorname{card}I'_{\alpha_0}\leq(\operatorname{card}\alpha_0)
\underbrace{\pi\leq\pi^2}_{\substack{\text{si }\operatorname{card}\alpha_0\leq\pi\text{ i.e.}\\\text{si }\pi_0<\pi}}=\pi,
\]
donc on gagne encore.

Ainsi, on a trouvé un deuxième système inductif
\[
(X_{\widetilde i})_{\widetilde i\in\widetilde I}
\]
de $=$ type que le système initial, ayant comme limite inductive $X$ (qu'il ait $=$ limite inductive est immédiat). Mais qu'a-t-on gagné? Ceci: \unsure{Dans|\deleted{Dans}} \deleted{$\unsure{I|X}$} $\widetilde I$ \unsure{tte|toute} \deleted{\unsure{Toute partie ordonnée|Toute partie}} partie \inserted{de $\widetilde I$ et \unsure{grande|grand} devant $\pi_0$} de cardinal $\leq\pi$ est non seulement \uline{majorée} dans $\widetilde I$, mais même y admet une borne supérieure.
\margin{$\Vert$}
\deleted{En effet, si on a des $I'_\alpha\in\widetilde I$, i.e. des $I'_\alpha\subset I$ qui}
\inserted{Cela résulte des}
